% =====================================================================
% MC714 - Sistemas Distribuídos - 2º Semestre de 2026
% Relatório 2 - Sistema Distribuído de Gerenciamento de Tarefas
% (Bully + Heartbeat + Primary-Backup)
%
% Arquivo "cru": apenas a estrutura de seções/subseções com comentários.
% Preâmbulo, classe IEEEtran, título e autores ficam na configuração
% já existente.
%
% LEMBRETES DO ENUNCIADO:
%  - Máximo de 6 páginas, duas colunas, SEM contar as referências.
%  - A PRIMEIRA PÁGINA deve conter:
%      (1) nome e RA dos participantes;
%      (2) link para o repositório no GitLab da disciplina;
%      (3) link para o vídeo de demonstração (~5 min, os dois integrantes).
%    Sugestão: colocar os links no \thanks{} do autor (como no Relatório 1)
%    ou em uma linha logo abaixo do abstract.
%  - É obrigatório citar as fontes de qualquer código reutilizado e
%    descrever as modificações feitas (ver Seção II-E).
%  - Pesos: correção e testes 30%; modelo, arquitetura e falhas 20%;
%    experimentos e reprodutibilidade 25%; relatório 15%; vídeo 10%.
%    Distribuir o espaço do texto de forma proporcional a esses pesos.
%
% Orçamento sugerido de páginas (aprox.):
%   I-II: 0,75 | III: 0,5 | IV-VI: 1,5 | VII: 0,5 | VIII: 0,5
%   IX-X: 1,75 | XI-XII: 0,5
% =====================================================================

% ---------------------------------------------------------------------
\begin{abstract}
% Um parágrafo (150-200 palavras), no mesmo formato do Relatório 1:
%  - O que foi construído: sistema de gerenciamento de tarefas com três
%    (ou mais) nós em contêineres Docker, comunicação HTTP (FastAPI/HTTPX)
%    e persistência SQLite por nó.
%  - Algoritmos: Bully (eleição), Primary-Backup síncrono (replicação) e
%    Heartbeat (detecção de falhas, bônus como terceiro algoritmo).
%  - Modelo adotado: parcialmente síncrono, crash-stop, uma falha por
%    ensaio, sem partições.
%  - Como foi avaliado: testes automatizados de correção com injeção de
%    falhas e experimentos variando N em {3,5,7} e o timeout de detecção.
%  - Um ou dois resultados numéricos principais (ex.: tempo médio de
%    eleição, aumento do p95 de escrita com N, tempo de detecção).
\end{abstract}

% Opcional (padrão IEEE):
% \begin{IEEEkeywords}
% Eleição de líder, Bully, replicação primário-backup, detecção de falhas,
% heartbeat, tolerância a falhas.
% \end{IEEEkeywords}

% =====================================================================
\section{Introdução}
% =====================================================================
% - Motivação: por que replicar um serviço e trocar de líder
%   automaticamente (disponibilidade x consistência).
% - Problema prático escolhido: API de tarefas (criar, listar, atualizar)
%   que deve continuar funcionando após a falha do servidor primário sem perder operações já confirmadas.
% - Quais algoritmos foram escolhidos e em quais categorias do enunciado
%   se encaixam: Bully -> eleição; Primary-Backup -> replicação;
%   Heartbeat -> detecção de falhas (terceiro algoritmo, bônus).
%   Mencionar também que a integração dos três é um item de bônus.
% - Contribuições em lista curta (implementação própria dos três
%   algoritmos, sem bibliotecas que já os forneçam; testes de falha;
%   experimentos reproduzíveis).
% - Parágrafo final com a organização do relatório (como no Relatório 1).

% =====================================================================
\section{Descrição da arquitetura do sistema}
% =====================================================================
% Parágrafo introdutório: visão geral em uma frase e referência à figura
% da arquitetura.

\subsection{Organização dos nós}
% - Figura (diagrama) com cliente, primário/líder e backups, cada um com
%   seu SQLite. Pode ser TikZ ou imagem exportada.
% - Todos os nós executam o mesmo código; mudam apenas NODE_ID, PEERS e
%   DB_PATH (variáveis de ambiente do compose.yaml).
% - Tabela pequena: nó, identificador, endereço interno, porta externa.
% - Papéis possíveis de um nó (ex.: PRIMARY, BACKUP, CANDIDATE/ELECTING,
%   SYNCING) e, se couber, um diagrama de estados simplificado.

\subsection{Componentes e responsabilidades}
% - Tabela: API HTTP, Bully, Heartbeat, Primary-Backup, SQLite,
%   gerenciador de estado (o papel de cada um).
% - Como os módulos interagem (Heartbeat aciona Bully; Bully dispara a
%   sincronização; Primary-Backup só é ativado depois dela).
% - Rotinas assíncronas que rodam em cada nó (asyncio): envio de
%   heartbeats, monitoramento, atendimento HTTP.

\subsection{Interface de comunicação}
% - Endpoints externos (cliente): POST /tasks, GET /tasks, PUT /tasks/{id};
%   comportamento quando o cliente acessa um backup (redireciona ou
%   informa o endereço do líder).
% - Endpoints internos: /health, /status, /internal/heartbeat,
%   /internal/election, /internal/coordinator, endpoints de PREPARE,
%   COMMIT e sincronização.
% - Formato das mensagens (JSON), campos comuns (node_id, epoch,
%   request_id, sequence_number).
% - Timeouts explícitos configurados no cliente HTTPX (o enunciado exige
%   deadlines/timeouts explícitos nas chamadas).

\subsection{Persistência}
% - Esquema SQLite: tabela de tarefas e tabela de operações (request_id,
%   epoch, sequence_number, operation_type, payload, status).
% - Estados da operação: PREPARED e COMMITTED.
% - Uso de transações para mudanças de status e atualização das tarefas.
% - O que é persistido sobre a configuração (epoch, líder, membros).

\subsection{Tecnologias e código reutilizado}
% - Linguagem e bibliotecas com versões: Python, FastAPI, Uvicorn, HTTPX,
%   SQLite, Docker Compose v2, Pandas/Matplotlib.
% - Deixar claro que nenhuma biblioteca utilizada implementa eleição,
%   replicação ou detecção de falhas (exigência do enunciado).
% - OBRIGATÓRIO: citar a fonte de qualquer código reutilizado (tutoriais,
%   exemplos, trechos gerados por ferramentas) e descrever as
%   modificações. Se não houver, declarar explicitamente.

% =====================================================================
\section{Modelo do sistema}
% =====================================================================
% O enunciado exige que o modelo seja declarado e JUSTIFICADO para cada
% algoritmo. Sugestão: descrever o modelo comum aqui e, nas seções de
% cada algoritmo, apenas as diferenças/hipóteses específicas.
% Uma tabela-resumo (característica x modelo adotado) economiza espaço.

\subsection{Modelo temporal}
% - Parcialmente síncrono: existem limites para atrasos e processamento,
%   mas não conhecidos a priori; os timeouts assumem que, nos ensaios em
%   que se espera correção, os atrasos ficam abaixo deles.
% - Justificar por que o Bully e o Heartbeat precisam dessa hipótese
%   (detecção por timeout).

\subsection{Propriedades dos canais}
% - HTTP sobre TCP: ordenação e ausência de duplicação dentro de uma
%   conexão; perdas tratadas como falha da requisição/timeout.
% - Duplicações no nível da aplicação (retentativas do cliente) tratadas
%   pelo request_id / Idempotency-Key.
% - Rede local entre contêineres (bridge do Docker).

\subsection{Modelo de falhas}
% - Crash-stop; no máximo uma falha de nó por ensaio.
% - Partições de rede NÃO consideradas; falhas bizantinas não consideradas.
% - Como as falhas são injetadas (docker compose stop/kill, falhas no
%   nível da aplicação, eventual uso de tc/netem ou Toxiproxy para
%   atrasos) — o enunciado pede que a escolha seja documentada.

\subsection{Relógios, persistência e reinicialização}
% - Relógio monotônico local apenas para timeouts; nenhuma sincronização
%   de relógios entre nós.
% - SQLite individual por nó (volume Docker).
% - Nós que falharam não são reintegrados (fora do escopo) — ou descrever
%   a reintegração, se implementada.

\subsection{Propriedades de segurança e vivacidade}
% Listar formalmente, por algoritmo (podem ser itens numerados S1, S2,
% L1... para serem referenciados nos testes):
% - Bully: segurança = no máximo um líder reconhecido por epoch;
%   vivacidade = após a falha do líder, um novo líder é eleito em tempo
%   finito (sob as hipóteses temporais).
% - Heartbeat: completude (todo nó que falha é eventualmente suspeito);
%   precisão (nenhum nó correto é suspeito, sob as hipóteses temporais).
% - Primary-Backup: nenhuma operação confirmada ao cliente é perdida;
%   réplicas com mesmo estado após sincronização; cada request_id aplicado
%   no máximo uma vez; mensagens de epochs antigos rejeitadas;
%   vivacidade = escritas voltam a ser aceitas após a reconfiguração.
% - Ressalva explícita: Bully + Heartbeat não constituem um protocolo de
%   consenso; não há garantia contra split-brain com partições ou falsas
%   suspeitas.

% =====================================================================
\section{Algoritmo de eleição: Bully}
% =====================================================================

\subsection{Funcionamento}
% - Ideia do algoritmo: o nó ativo de maior ID torna-se coordenador.
% - Mensagens ELECTION, OK e COORDINATOR (tabela curta).
% - Diagrama de sequência do exemplo (falha do Node 3, Node 1 inicia,
%   Node 2 vence).

\subsection{Implementação}
% - Pseudocódigo (algorithm/algorithmic ou listing) da rotina de eleição.
% - OK retornado como resposta HTTP ao /internal/election.
% - Tratamento de: eleições simultâneas, recebimento de ELECTION durante
%   uma eleição, timeout de espera pelo COORDINATOR e reinício da
%   eleição, COORDINATOR atrasado ou de nó com ID menor.
% - Incremento do epoch ao anunciar a liderança.
% - Ponto importante: o líder eleito só se torna PRIMARY após a
%   sincronização (Seção VII).

\subsection{Complexidade}
% - Mensagens: O(N) no melhor caso (maior ID detecta a falha) e O(N^2)
%   no pior caso (menor ID inicia). Escrever a contagem esperada para
%   N = 3, 5, 7, que será comparada com os experimentos.
% - Tempo: dominado pelos timeouts de espera por OK/COORDINATOR.

% =====================================================================
\section{Algoritmo de detecção de falhas: Heartbeat}
% =====================================================================

\subsection{Funcionamento}
% - Envio periódico (todos para todos) e registro do instante da última
%   mensagem recebida de cada participante.
% - Suspeita quando o tempo desde o último heartbeat excede o timeout.
% - Heartbeat como mecanismo independente que apenas aciona o Bully.

\subsection{Implementação e parâmetros}
% - HEARTBEAT_INTERVAL e FAILURE_TIMEOUT (valores padrão e justificativa).
% - Envios concorrentes com timeout, para que um nó indisponível não
%   atrase os demais.
% - Como se evita processar a mesma falha repetidamente.
% - Custo: N(N-1) mensagens por intervalo.

\subsection{Integração com o Bully}
% - Falha do líder -> marcar como suspeito -> iniciar eleição.
% - Falha de um backup -> sem eleição; líder inicia reconfiguração
%   (remove o nó dos membros ativos após o procedimento da Seção VII).
% - Discussão: o Heartbeat produz suspeitas, não provas; relação entre
%   timeout curto e falsas suspeitas.

% =====================================================================
\section{Algoritmo de replicação: Primary-Backup}
% =====================================================================

\subsection{Modelo de dados e operações}
% - Representação de uma tarefa (id, title, status, version).
% - Operações da API e quais são replicadas (criação e atualização).
% - Apenas o líder atende leituras e escritas (justificar: simplifica a
%   consistência; leituras sempre refletem escritas confirmadas).

\subsection{Replicação síncrona em duas fases}
% - Figura do fluxo PREPARE -> ACK -> COMMIT -> COMMIT ACK -> HTTP 201.
% - Fase A (preparação) e Fase B (confirmação), passo a passo.
% - Confirmação por TODOS os membros da configuração ativa; justificar a
%   escolha (garantia simples diante de uma única falha).
% - Deixar explícito que não é uma implementação completa de 2PC: as fases
%   fazem parte do mecanismo de replicação.
% - Comportamento quando um nó falha no meio da operação (erro ou
%   resultado desconhecido; nunca sucesso sem as confirmações).

\subsection{Mecanismos contra inconsistências}
% - Serialização das escritas (asyncio.Lock) e sequence_number.
% - Identificação única (request_id / Idempotency-Key) e reenvios seguros.
% - Epoch em toda mensagem de replicação; backups rejeitam epochs antigos
%   e líderes não reconhecidos.

\subsection{Complexidade}
% - Mensagens por escrita: 4(N-1) (PREPARE, ACK, COMMIT, COMMIT ACK).
% - Latência dominada pelo backup mais lento (espera por todos).
% - Número de escritas em disco por operação.

% =====================================================================
\section{Integração e recuperação após falhas}
% =====================================================================
% Seção que conecta os três algoritmos — é o cenário central do trabalho
% e um item de bônus ("integração para resolver um problema prático").

\subsection{Gerenciamento de configurações}
% - Configuração explícita {epoch, leader, members}; exemplo antes e
%   depois da falha do Node 3.
% - Suspensão das escritas durante a transição.
% - Rejeição de mensagens de configurações antigas.

\subsection{Recuperação após a falha do primário}
% - Passo a passo: falha -> detecção -> eleição -> sincronização ->
%   ativação -> continuidade.
% - Sincronização: consulta aos logs dos sobreviventes, identificação das
%   operações COMMITTED, resolução das PREPARED (recuperar ou descartar,
%   com justificativa), atualização do backup e confirmação da nova
%   configuração.
% - Argumento de correção: por que nenhuma operação confirmada ao cliente
%   é perdida sob o modelo de uma única falha.

\subsection{Falha de um backup}
% - Detecção pelo líder, escrita em andamento não é confirmada, nova
%   configuração sem o backup, retomada das escritas.

\subsection{Limites das garantias}
% - O que acontece fora do modelo (partição, falsa suspeita, duas falhas):
%   possibilidade de dois líderes, perda de disponibilidade etc.
% - Comparação breve com protocolos de consenso com quóruns (Raft/Paxos).

% =====================================================================
\section{Testes de correção}
% =====================================================================

\subsection{Metodologia}
% - Testes unitários (lógica de cada algoritmo) e de integração (cluster
%   Docker real).
% - Script de injeção de falhas: sequência iniciar -> aguardar eleição ->
%   operações -> derrubar nó -> aguardar -> operações -> verificar ->
%   encerrar.
% - Verificações automáticas via /status e comparação dos bancos de
%   cada nó; timeouts de teste para detectar não convergência.
% - Logs estruturados e como são usados como evidência (o enunciado
%   aceita verificação reproduzível dos logs).

\subsection{Cenários e resultados}
% - Tabela: cenário | procedimento | resultado esperado | resultado obtido
%   (ou nº de execuções aprovadas / total). Cenários:
%     eleição inicial; falha do líder; replicação; continuidade
%     (100 + 100 tarefas); falha entre PREPARE e COMMIT; falha de backup
%     durante operação; requisição duplicada; eleições simultâneas;
%     COORDINATOR atrasado.
% - Associar cada cenário às propriedades S/L definidas na Seção III-E.
% - O enunciado exige, por algoritmo, pelo menos um cenário normal e um de
%   falha — conferir que a tabela cobre isso para os três.
% - Comentar falhas encontradas durante o desenvolvimento e como foram
%   corrigidas (se houver espaço).

% =====================================================================
\section{Metodologia experimental}
% =====================================================================

\subsection{Ambiente de execução}
% - Hardware (CPU, núcleos, RAM), sistema operacional, versões de Docker,
%   Docker Compose, Python e bibliotecas.
% - Observação obrigatória: contêineres em uma única máquina são nós
%   lógicos, não domínios de falha independentes; latência de rede
%   artificialmente baixa.
% - Se usada, comparação com a nuvem privada do IC (bônus).

\subsection{Automação e reprodutibilidade}
% - Scripts (run_experiments.py, inject_failures.py, analyze_results.py)
%   e comando único para reproduzir.
% - Número de repetições por configuração (>= 10), sementes registradas,
%   duração dos ensaios e período de aquecimento.
% - Onde ficam os dados brutos no repositório e a tag de entrega.

\subsection{Métricas e parâmetros}
% - Tabela: algoritmo | parâmetro variado (3 níveis) | parâmetros fixos |
%   métricas.
%     Bully: N em {3,5,7}; tempo de eleição, nº de mensagens por tipo,
%            taxa de eleições corretas.
%     Primary-Backup: N em {3,5,7}, carga de 1000 escritas; latência média,
%            p50/p95/p99, throughput, mensagens por operação.
%     Heartbeat: T em {1,5; 2,5; 4,0} s com intervalo de 500 ms; tempo de
%            detecção (média e percentis), taxa de falsas suspeitas.
% - Justificar cada métrica (por que ela caracteriza o algoritmo).
% - Definir precisamente como cada tempo é medido (ex.: tempo de eleição
%   = instante do stop até todos os sobreviventes reconhecerem o novo
%   líder; quais relógios são usados).
% - Medidas de variabilidade: desvio padrão ou IC de 95% (como no
%   Relatório 1).

% =====================================================================
\section{Resultados experimentais}
% =====================================================================
% Para cada subseção: gráfico(s) + tabela compacta + análise comparando
% com o esperado e com a complexidade teórica (Seções IV-D, V-B e VI-D).

\subsection{Bully}
% - Gráficos: tempo de eleição x N; mensagens x N (separando ELECTION,
%   OK, COORDINATOR) com a curva teórica sobreposta.
% - Discutir o crescimento observado e o efeito de quem inicia a eleição.

\subsection{Primary-Backup}
% - Gráficos: latência (média, p50, p95, p99) x N; throughput x N.
% - Tabela com os percentis.
% - Discutir o impacto da espera pelo backup mais lento e das escritas
%   síncronas em disco.

\subsection{Heartbeat}
% - Gráficos: tempo de detecção x timeout (com percentis); falsas
%   suspeitas x timeout.
% - Comparar o tempo de detecção medido com o esperado
%   (aprox. entre T e T + intervalo de verificação).
% - Discutir o compromisso rapidez x falsas suspeitas.

\subsection{Cenário integrado}
% - (Opcional, se couber) tempo total de indisponibilidade das escritas
%   após a falha do primário, decomposto em detecção + eleição +
%   sincronização.

% =====================================================================
\section{Discussão e limitações}
% =====================================================================
% - Bully + Heartbeat não são consenso: possibilidade de split-brain com
%   partições ou falsas suspeitas.
% - Replicação para todos os membros: uma falha de backup bloqueia as
%   escritas até a reconfiguração (disponibilidade x consistência).
% - Primário único como gargalo de escrita e leitura.
% - Sem reintegração de nós que falharam (se for o caso).
% - Limitações do ambiente (uma máquina, rede local, ruído de medição).
% - Possíveis extensões: quóruns, reintegração com transferência de
%   estado, leituras em backups, Raft.

% =====================================================================
\section{Conclusão}
% =====================================================================
% - Retomar o objetivo e o que foi demonstrado (troca de líder automática
%   sem perda de operações confirmadas, dentro do modelo).
% - Principais resultados quantitativos de cada algoritmo em uma frase.
% - Lições aprendidas e trabalhos futuros (curto).

% =====================================================================
% Referências (não contam no limite de 6 páginas)
% =====================================================================
% Sugestões:
%  - M. van Steen e A. S. Tanenbaum, Distributed Systems, 4th ed., 2023.
%  - H. Garcia-Molina, "Elections in a Distributed Computing System",
%    IEEE Trans. Computers, 1982 (artigo original do Bully).
%  - N. Budhiraja et al., "The Primary-Backup Approach", in Distributed
%    Systems (Mullender, ed.), 1993.
%  - T. D. Chandra e S. Toueg, "Unreliable Failure Detectors for
%    Reliable Distributed Systems", JACM, 1996.
%  - Documentação do FastAPI, HTTPX, SQLite e Docker Compose.
%  - Fontes de qualquer código reutilizado.
%
% \begin{thebibliography}{00}
% \bibitem{b1} ...
% \end{thebibliography}